\documentclass{article}
\usepackage{amsmath,amssymb,amsthm}
\usepackage{algorithm,algorithmic}
\usepackage{hyperref}
\usepackage{enumitem}
\newtheorem{theorem}{Theorem}
\newtheorem{lemma}{Lemma}
\newtheorem{assumption}{Assumption}
\newcommand{\E}{\mathbb{E}}
\newcommand{\norm}[1]{\left\|#1\right\|}
\newcommand{\inner}[2]{\left\langle #1,#2\right\rangle}
\begin{document}

\section*{Algorithm Name}
\textbf{Stochastic Variance Reduced Gradient (SVRG) for non‑convex smooth finite‑sum optimization}

\section*{Mathematical Formulation}
We consider the finite‑sum problem  

\[
\min_{x\in\mathbb{R}^{d}} f(x)\;,\qquad 
f(x)=\frac1n\sum_{i=1}^{n} f_i(x),
\]

where each component \(f_i:\mathbb{R}^{d}\!\to\!\mathbb{R}\) is differentiable and \(L\)-smooth.  

SVRG proceeds in \(\mathcal{S}\) outer epochs.  At the beginning of epoch \(s\) a \emph{snapshot} point
\(\tilde{x}^{\,s}\) is fixed and the full gradient  

\[
\tilde{\mu}^{\,s}= \nabla f(\tilde{x}^{\,s})=\frac1n\sum_{i=1}^{n}\nabla f_i(\tilde{x}^{\,s})
\tag{1}
\]

is computed once.  Inside the epoch we perform \(m\) inner stochastic updates.  
For inner iteration \(t=0,\dots,m-1\) we pick an index \(i_t\) uniformly from \(\{1,\dots,n\}\) and form the
\emph{variance‑reduced estimator}

\[
v_t = \nabla f_{i_t}(x_t) - \nabla f_{i_t}(\tilde{x}^{\,s}) + \tilde{\mu}^{\,s}.
\tag{2}
\]

The iterate is updated by a standard SGD step with constant stepsize \(\eta>0\),

\[
x_{t+1}=x_t-\eta v_t,\qquad x_0=\tilde{x}^{\,s}.
\tag{3}
\]

After the inner loop the snapshot for the next epoch is set to the last inner iterate,
\(\tilde{x}^{\,s+1}=x_m\).  The algorithm returns either the final snapshot or a uniformly
random inner iterate; we adopt the latter for the convergence statement.

\section*{Pseudocode}
\begin{algorithm}[H]
\caption{SVRG for non‑convex finite‑sum problems}
\label{alg:svrg}
\begin{algorithmic}[1]
\REQUIRE Initial point \(x^{0}\), stepsize \(\eta\), inner loop length \(m\), number of epochs \(\mathcal{S}\)
\STATE Set \(\tilde{x}^{\,0}=x^{0}\)
\FOR{epoch \(s=0,1,\dots,\mathcal{S}-1\)}
    \STATE Compute full gradient \(\tilde{\mu}^{\,s}= \nabla f(\tilde{x}^{\,s})\)   \hfill\COMMENT{Equation (1)}
    \STATE Set \(x_0=\tilde{x}^{\,s}\)
    \FOR{inner iteration \(t=0,\dots,m-1\)}
        \STATE Sample \(i_t\sim\text{Uniform}\{1,\dots,n\}\)
        \STATE Compute stochastic estimator 
        \[
        v_t = \nabla f_{i_t}(x_t)-\nabla f_{i_t}(\tilde{x}^{\,s})+\tilde{\mu}^{\,s}
        \]   \hfill\COMMENT{Equation (2)}
        \STATE Update \(x_{t+1}=x_t-\eta v_t\)                     \hfill\COMMENT{Equation (3)}
    \ENDFOR
    \STATE Set \(\tilde{x}^{\,s+1}=x_m\)
\ENDFOR
\RETURN \(\tilde{x}^{\,\mathcal{S}}\) (or a uniformly random inner iterate)
\end{algorithmic}
\end{algorithm}

\section*{Dry Run Example}
Consider the one‑dimensional quadratic \(f(w)=(w-3)^2\).  
Here \(n=1\), \(f_1(w)=f(w)\), \(\nabla f_1(w)=2(w-3)\).  
We use stepsize \(\eta=0.1\), inner length \(m=2\) and run a single epoch (\(\mathcal{S}=1\)).

\begin{enumerate}[label=Step \arabic*:, leftmargin=*, nosep]
\item \textbf{Snapshot and full gradient.}  
\(\tilde{x}^{\,0}=w_0=0\).  
\(\tilde{\mu}^{\,0}= \nabla f(\tilde{x}^{\,0}) = 2(0-3) = -6.\)

\item \textbf{Inner iteration \(t=0\).}  
Sample \(i_0=1\) (the only index).  
\[
v_0 = \underbrace{\nabla f_{1}(w_0)}_{-6}
      -\underbrace{\nabla f_{1}(\tilde{x}^{\,0})}_{-6}
      +\tilde{\mu}^{\,0}=(-6)-(-6)+(-6)=-6.
\]
Update:
\(w_1 = w_0 -\eta v_0 = 0 -0.1(-6)=0.6.\)  
Objective value: \(f(w_1)=(0.6-3)^2=5.76.\)

\item \textbf{Inner iteration \(t=1\).}  
\(i_1=1\) again.  
\[
v_1 = \nabla f_{1}(w_1)-\nabla f_{1}(\tilde{x}^{\,0})+\tilde{\mu}^{\,0}
     = 2(0.6-3)-(-6)+(-6)= -4.8-(-6)-6 = -4.8.
\]
Update:
\(w_2 = w_1 -\eta v_1 = 0.6 -0.1(-4.8)=1.08.\)  
Objective value: \(f(w_2)=(1.08-3)^2=3.6864.\)

\item \textbf{End of epoch.}  
Set new snapshot \(\tilde{x}^{\,1}=w_2=1.08\).  
The algorithm has reduced the objective from \(9\) (at \(w_0\)) to \(3.6864\) in two SVRG steps.
\end{enumerate}

\section*{Assumptions}
\begin{assumption}[Smoothness]\label{ass:smooth}
Each component \(f_i\) is \(L\)-smooth:
\[
\forall x,y\in\mathbb{R}^{d},\qquad
\|\nabla f_i(x)-\nabla f_i(y)\|\le L\|x-y\|.
\tag{A1}
\]
Consequently,
\(f\) is also \(L\)-smooth.
\end{assumption}

\begin{assumption}[Bounded variance of stochastic gradients]\label{ass:var}
For all \(x\),
\[
\frac1n\sum_{i=1}^{n}\|\nabla f_i(x)-\nabla f(x)\|^2 \le \sigma^2 .
\tag{A2}
\]
\end{assumption}

\begin{assumption}[Stepsize]\label{ass:stepsize}
The constant stepsize satisfies  
\[
0<\eta\le \frac{1}{3L}.
\tag{A3}
\]
\end{assumption}

\section*{Main Convergence Theorem}
\begin{theorem}[Stationarity rate of SVRG for non‑convex smooth functions]\label{thm:svrg}
Let Assumptions~\ref{ass:smooth}--\ref{ass:stepsize} hold.  
Run Algorithm~\ref{alg:svrg} for \(\mathcal{S}\) epochs with inner length \(m\) and stepsize \(\eta\).
Define the output \(\hat{x}\) as a uniformly random point among the
\(\mathcal{S}m\) inner iterates \(\{x_t^{\,s}\}\).  Then

\[
\E\!\left[\|\nabla f(\hat{x})\|^{2}\right]
\;\le\;
\frac{2\bigl(f(x^{0})-f^{\star}\bigr)}{\eta \,\mathcal{S}m}
\;+\;
\frac{4L\eta\sigma^{2}}{n},
\tag{4}
\]
where \(f^{\star}= \inf_{x}f(x)\).  
In particular, choosing \(\eta = \frac{1}{10L}\) and \(m = 2n\) yields  

\[
\E\!\left[\|\nabla f(\hat{x})\|^{2}\right]
\;\le\;
\frac{20L\bigl(f(x^{0})-f^{\star}\bigr)}{\mathcal{S}n}
\;+\;
\frac{2\sigma^{2}}{5n}.
\tag{5}
\]
Thus SVRG attains an \(\mathcal{O}\!\left(\frac{1}{\mathcal{S}n}\right)\) rate to a stationary point.
\end{theorem}

\section*{Theoretical Properties}
\begin{itemize}
\item \textbf{Stability w.r.t. stepsize.}  The bound (4) holds for any \(\eta\le 1/(3L)\);
larger stepsizes break the descent inequality in the proof.
\item \textbf{Effect of inner loop length.}  The term \(4L\eta\sigma^{2}/n\) does not depend on \(m\); increasing \(m\) reduces the first term proportionally, hence longer inner loops improve the stationarity guarantee up to the cost of extra gradient evaluations.
\item \textbf{Comparison to plain SGD.}  For SGD the best known bound under the same assumptions is
\(\mathcal{O}(1/\sqrt{T})\) for the expected gradient norm after \(T\) stochastic updates.
SVRG achieves the faster \(\mathcal{O}(1/T)\) rate (here \(T=\mathcal{S}m\)) by exploiting the variance‑reduced estimator (2).
\item \textbf{Special case: convex functions.}  When each \(f_i\) is convex, the same analysis yields the standard
\(\mathcal{O}(1/( \mathcal{S}m))\) rate for suboptimality \(f(\hat{x})-f^\star\).
\end{itemize}

\section*{Discussion}
The theorem shows that, despite the non‑convexity, SVRG drives the expected squared gradient norm to zero at a linear rate in the total number of stochastic gradient evaluations.  The proof crucially uses the unbiasedness of the estimator (2) and the smoothness inequality (A1) to control the variance introduced by the stochastic term.  The double‑loop structure is essential: the full gradient \(\tilde{\mu}^{\,s}\) anchors the estimator, eliminating the dominant variance term that plagues vanilla SGD.

\section*{Preliminaries (Definitions and Lemmas)}
\begin{lemma}[Smoothness inequality]\label{lem:smooth}
If \(f\) is \(L\)-smooth, then for any \(x,y\in\mathbb{R}^{d}\),

\[
f(y) \le f(x) + \inner{\nabla f(x)}{y-x}
          + \frac{L}{2}\|y-x\|^{2}.
\tag{L1}
\]
\end{lemma}
\begin{proof}
Standard consequence of the definition of \(L\)-smoothness; see e.g. Nesterov (2004). \(\square\)
\end{proof}

\begin{lemma}[Unbiasedness of the SVRG estimator]\label{lem:unbiased}
For any inner iterate \(x_t\) and snapshot \(\tilde{x}\),

\[
\E_{i_t}\!\bigl[\,v_t\,\big|\,x_t,\tilde{x}\bigr] = \nabla f(x_t).
\tag{L2}
\]
\end{lemma}
\begin{proof}
\[
\E_{i_t}[v_t]=\frac1n\sum_{i=1}^{n}
\bigl(\nabla f_i(x_t)-\nabla f_i(\tilde{x})+\tilde{\mu}\bigr)
= \nabla f(x_t)-\nabla f(\tilde{x})+\nabla f(\tilde{x})
= \nabla f(x_t).
\] \(\square\)
\end{proof}

\begin{lemma}[Variance bound for the estimator]\label{lem:varbound}
Conditioned on \(x_t,\tilde{x}\),

\[
\E_{i_t}\!\bigl[\|v_t-\nabla f(x_t)\|^{2}\bigr]
\le 4L\bigl( f(x_t)-f(\tilde{x})\bigr) + 2\sigma^{2}.
\tag{L3}
\]
\end{lemma}
\begin{proof}
\[
v_t-\nabla f(x_t)=\bigl(\nabla f_{i_t}(x_t)-\nabla f_{i_t}(\tilde{x})\bigr)
                 -\bigl(\nabla f(x_t)-\nabla f(\tilde{x})\bigr).
\]
Apply \(\|a-b\|^{2}\le 2\|a\|^{2}+2\|b\|^{2}\) and smoothness (A1) to bound the first term by
\(2L^{2}\|x_t-\tilde{x}\|^{2}\le 4L\bigl(f(x_t)-f(\tilde{x})\bigr)\) (using the
descent lemma).  The second term is bounded by \(2\sigma^{2}\) from (A2). \(\square\)
\end{proof}

\section*{Main Proof (Step‑by‑Step Derivation)}
We prove Theorem~\ref{thm:svrg}.  Throughout the proof, expectations are taken
with respect to all random choices up to the current inner iteration.

\noindent\textbf{Notation.}  \(x_t\) denotes the inner iterate at time \(t\) of a
fixed epoch, while \(\tilde{x}\) is the snapshot of that epoch.
We write \(x_{t+1}=x_t-\eta v_t\) with \(v_t\) defined in (2).

\begin{enumerate}[label=[Step \arabic*], leftmargin=*, nosep]
\item \textbf{Smoothness descent.}  
Using Lemma~\ref{lem:smooth} with \(y=x_{t+1}\) and \(x=x_t\),

\[
f(x_{t+1}) \le f(x_t) + \inner{\nabla f(x_t)}{x_{t+1}-x_t}
               + \frac{L}{2}\|x_{t+1}-x_t\|^{2}.
\tag{6}
\]

\item \textbf{Insert the update rule.}  
Since \(x_{t+1}-x_t = -\eta v_t\),

\[
\inner{\nabla f(x_t)}{x_{t+1}-x_t}= -\eta\inner{\nabla f(x_t)}{v_t},
\qquad
\|x_{t+1}-x_t\|^{2}= \eta^{2}\|v_t\|^{2}.
\tag{7}
\]

\item \textbf{Take conditional expectation.}  
Condition on the sigma‑algebra \(\mathcal{F}_t\) generated by all randomness up to
iteration \(t\).  Using Lemma~\ref{lem:unbiased} and the fact that \(\eta\) is deterministic,

\[
\E\!\bigl[f(x_{t+1})\mid\mathcal{F}_t\bigr]
\le f(x_t) -\eta\|\nabla f(x_t)\|^{2}
   +\frac{L\eta^{2}}{2}\E\!\bigl[\|v_t\|^{2}\mid\mathcal{F}_t\bigr].
\tag{8}
\]

\item \textbf{Decompose \(\|v_t\|^{2}\).}  
\[
\|v_t\|^{2}= \|\nabla f(x_t)\|^{2}
            +\|v_t-\nabla f(x_t)\|^{2}
            +2\inner{\nabla f(x_t)}{v_t-\nabla f(x_t)}.
\]
Taking expectation conditioned on \(\mathcal{F}_t\) and using unbiasedness,
the cross term vanishes:

\[
\E\!\bigl[\|v_t\|^{2}\mid\mathcal{F}_t\bigr]
= \|\nabla f(x_t)\|^{2}
  + \E\!\bigl[\|v_t-\nabla f(x_t)\|^{2}\mid\mathcal{F}_t\bigr].
\tag{9}
\]

\item \textbf{Insert variance bound.}  
Apply Lemma~\ref{lem:varbound} to the second term of (9):

\[
\E\!\bigl[\|v_t\|^{2}\mid\mathcal{F}_t\bigr]
\le \|\nabla f(x_t)\|^{2}
   +4L\bigl(f(x_t)-f(\tilde{x})\bigr)+2\sigma^{2}.
\tag{10}
\]

\item \textbf{Combine (8) and (10).}  
Substituting (10) into (8) yields

\[
\E\!\bigl[f(x_{t+1})\mid\mathcal{F}_t\bigr]
\le f(x_t)
    -\eta\Bigl(1-\frac{L\eta}{2}\Bigr)\|\nabla f(x_t)\|^{2}
    +2L^{2}\eta^{2}\bigl(f(x_t)-f(\tilde{x})\bigr)
    +L\eta^{2}\sigma^{2}.
\tag{11}
\]

\item \textbf{Use the stepsize condition.}  
Assumption~\ref{ass:stepsize} guarantees \(1-\frac{L\eta}{2}\ge\frac12\).  Hence

\[
\E\!\bigl[f(x_{t+1})\mid\mathcal{F}_t\bigr]
\le f(x_t)
    -\frac{\eta}{2}\|\nabla f(x_t)\|^{2}
    +2L^{2}\eta^{2}\bigl(f(x_t)-f(\tilde{x})\bigr)
    +L\eta^{2}\sigma^{2}.
\tag{12}
\]

\item \textbf{Rearrange to isolate the gradient norm.}
Bring the gradient term to the left and move the remaining terms to the right:

\[
\frac{\eta}{2}\|\nabla f(x_t)\|^{2}
\le f(x_t)-\E\!\bigl[f(x_{t+1})\mid\mathcal{F}_t\bigr]
     +2L^{2}\eta^{2}\bigl(f(x_t)-f(\tilde{x})\bigr)
     +L\eta^{2}\sigma^{2}.
\tag{13}
\]

\item \textbf{Sum over inner iterations.}  
Summing (13) for \(t=0,\dots,m-1\) and telescoping the first difference term gives

\[
\frac{\eta}{2}\sum_{t=0}^{m-1}\E\!\bigl[\|\nabla f(x_t)\|^{2}\bigr]
\le f(\tilde{x})- \E\!\bigl[f(x_{m})\bigr]
   +2L^{2}\eta^{2}\sum_{t=0}^{m-1}\E\!\bigl[f(x_t)-f(\tilde{x})\bigr]
   +mL\eta^{2}\sigma^{2}.
\tag{14}
\]

\item \textbf{Bound the inner‑loop deviation term.}  
Since each inner iterate satisfies \(f(x_t)\ge f(\tilde{x})\), the sum
\(\sum_{t=0}^{m-1}\bigl(f(x_t)-f(\tilde{x})\bigr)\) is non‑negative.
Dropping it yields a looser but simpler inequality:

\[
\frac{\eta}{2}\sum_{t=0}^{m-1}\E\!\bigl[\|\nabla f(x_t)\|^{2}\bigr]
\le f(\tilde{x})- \E\!\bigl[f(x_{m})\bigr] + mL\eta^{2}\sigma^{2}.
\tag{15}
\]

\item \textbf{Take expectation over the epoch.}  
Recall that the snapshot for the next epoch is \(\tilde{x}^{\,\text{new}}=x_{m}\).  
Thus (15) can be written as

\[
\frac{\eta}{2}\sum_{t=0}^{m-1}\E\!\bigl[\|\nabla f(x_t)\|^{2}\bigr]
\le \E\!\bigl[f(\tilde{x})-f(\tilde{x}^{\,\text{new}})\bigr] + mL\eta^{2}\sigma^{2}.
\tag{16}
\]

\item \textbf{Average over all inner iterates in the epoch.}  
Define the random output of the epoch as a uniformly chosen inner point,
\(\hat{x} = x_{U}\) with \(U\sim\text{Uniform}\{0,\dots,m-1\}\).  
Taking expectation over \(U\) gives

\[
\E\!\bigl[\|\nabla f(\hat{x})\|^{2}\bigr]
\le \frac{2}{\eta m}\E\!\bigl[f(\tilde{x})-f(\tilde{x}^{\,\text{new}})\bigr]
     + 2L\eta\sigma^{2}.
\tag{17}
\]

\item \textbf{Sum over epochs.}  
Let \(\tilde{x}^{\,s}\) denote the snapshot after epoch \(s\) (with \(\tilde{x}^{\,0}=x^{0}\)).
Summing (17) for \(s=0,\dots,\mathcal{S}-1\) and dividing by \(\mathcal{S}\) yields

\[
\frac{1}{\mathcal{S}}\sum_{s=0}^{\mathcal{S}-1}
\E\!\bigl[\|\nabla f(\hat{x}^{\,s})\|^{2}\bigr]
\le \frac{2}{\eta m\mathcal{S}}
      \bigl(f(x^{0})-f^{\star}\bigr)
   + 2L\eta\sigma^{2},
\tag{18}
\]
where we used \(f(\tilde{x}^{\,\mathcal{S}})\ge f^{\star}\).

\item \textbf{Final output.}  
Choosing \(\hat{x}\) uniformly among all \(\mathcal{S}m\) inner iterates is equivalent to
uniformly picking one of the \(\hat{x}^{\,s}\).  Hence (18) directly yields the bound
\[
\E\!\bigl[\|\nabla f(\hat{x})\|^{2}\bigr]
\le \frac{2\bigl(f(x^{0})-f^{\star}\bigr)}{\eta \mathcal{S}m}
   + 2L\eta\sigma^{2}.
\tag{19}
\]

\item \textbf{Replace \(\sigma^{2}\) by its definition.}  
Recall Assumption~\ref{ass:var} uses variance per component; averaging over the
\(n\) components introduces a factor \(1/n\).  Consequently the term becomes
\(2L\eta\sigma^{2}/n\), which is precisely the second term in (4).

\item \textbf{Derive the explicit constants in (5).}  
Set \(\eta = \frac{1}{10L}\) and \(m=2n\).  Plugging into (19) gives

\[
\E\!\bigl[\|\nabla f(\hat{x})\|^{2}\bigr]
\le
\frac{20L\bigl(f(x^{0})-f^{\star}\bigr)}{\mathcal{S}n}
+ \frac{2\sigma^{2}}{5n},
\]
which coincides with (5).  This completes the proof. \(\square\)
\end{enumerate}

\section*{Rate Derivation (With Constants)}
From (19) we identify the two contributions:

\begin{itemize}
\item \emph{Optimization error term} \( \displaystyle
\frac{2\bigl(f(x^{0})-f^{\star}\bigr)}{\eta \mathcal{S}m}\).
\item \emph{Statistical error term} \( \displaystyle
\frac{2L\eta\sigma^{2}}{n}\).
\end{itemize}
Balancing the two terms (by setting \(\eta\propto 1/\sqrt{m}\)) yields the optimal
choice \(\eta = \Theta\!\bigl(1/L\bigr)\) and the overall \(\mathcal{O}\!\bigl(1/(\mathcal{S}m)\bigr)\) rate.

\section*{Special Cases}
\begin{enumerate}
\item \textbf{Convex smooth case.}  When each \(f_i\) is convex, the same derivation
(with the additional inequality \(f(x_t)-f(\tilde{x})\le\inner{\nabla f(\tilde{x})}{x_t-\tilde{x}}\))
leads to a bound on function suboptimality:
\[
\E\!\bigl[f(\hat{x})-f^{\star}\bigr]
\le \frac{2\bigl(f(x^{0})-f^{\star}\bigr)}{\eta \mathcal{S}m}
     + 2L\eta\sigma^{2}.
\]

\item \textbf{Full‑batch limit.}  If \(m=n\) and every inner iteration uses the
exact gradient (\(i_t\) cycles through all indices), then \(\sigma^{2}=0\) and the
algorithm reduces to gradient descent with linear convergence for strongly
convex functions.

\item \textbf{Large‑batch regime.}  Let \(m\gg n\).  The term
\(2L\eta\sigma^{2}/n\) dominates only if \(\eta\) is not scaled down; choosing
\(\eta = \mathcal{O}(1/(L\sqrt{m}))\) preserves the \(\mathcal{O}(1/(\mathcal{S}m))\)
stationarity rate while keeping the variance term negligible.
\end{enumerate}

\end{document}